% ============================================================
%  chapters/ch03_methodology.tex — Chapter 3: Methodology
% ============================================================

\chapter{Methodology}
% ------------------------------------------------------------
\section{Research Design Overview}
This study employs an upstream computational audit design centered on stateless API execution to quantify model-level system bias across commercial foundation models.

\section{Task-Anchored Matrix Operationalization}
Adopting \cite{Jain2025} task taxonomy to operationalize Graduate Teaching Assistant (GTA) workflows across two distinct problem spaces:

\subsection{Task 1 --- Objective Context}
\begin{itemize}
  \item \textbf{Category D:} Problem-Solving Objective.
  \item \textbf{Scenario:} Routine Concept Explanation (e.g., explaining the K-Nearest Neighbors [KNN] algorithm in an introductory Data Science course).
  \item \textbf{Target State:} Fixed core ground truth. Diversity is defined strictly by pedagogical explanatory variety (diverse analogies and framings) while maintaining core domain accuracy without conceptual distortion.
\end{itemize}

\subsection{Task 2 --- Subjective Context}
\begin{itemize}
  \item \textbf{Category E:} Problem-Solving Subjective.
  \item \textbf{Scenario:} Open-Ended Research Scoping (e.g., guiding semester project ideation and research question formulation for an introductory AI Ethics course).
  \item \textbf{Target State:} Open, value-laden solution space. Diversity is defined by perspective pluralism (generating distinct, personalized analytical angles) versus epistemic enclosure (collapsing student ideation into a single, statistically dominant narrative).
\end{itemize}

\section{Model and Experimental Setup}

\subsection{Foundation Model Pool and Selection Rationale}
Leading commercial LLMs representing dominant technological baselines were selected: OpenAI GPT-4o, Anthropic Claude 3.5 Sonnet, and Google Gemini 1.5 Pro.

\subsection{Prompt Standardization Protocol}
Adopting \cite{Sahoo2024} Zero-Shot Task Framing Protocol to isolate model-level architectural bias from prompt engineering variation.

\subsection{Sampling Matrix and Session Isolation}
The experimental matrix comprises:
\begin{equation}
3 \text{ Models} \times 2 \text{ Tasks} \times 50 \text{ Iterations} = 300 \text{ Independent API Outputs}
\end{equation}
Stateless API execution with zero context memory was enforced to guarantee independent sample generation.

\subsection{Hyperparameter Standardization Protocol}
Decoding parameters (Temperature, $\text{Top\_p}$, and Max Tokens) were strictly controlled across all model endpoints to decouple generation stochasticity from systemic model bias, ensuring strict cross-architecture comparability.

\section{Data Preprocessing and Corpus Construction}
Preprocessing pipeline includes Markdown stripping, token normalization, and structured JSON vector corpus preparation.

\section{Quantitative and Computational Metrics}

\subsection{Measuring Homogenization and Diversity Collapse (RQ2/H2)}
High-dimensional Sentence Embedding Variance (Cosine Distance) combined with BERTopic Cluster Diversity quantifies semantic vector space shrinkage in Task 2.

\subsection{Measuring Explanatory Variety and Perspective Coverage (RQ1/H1)}
Quantitative Content Analysis (QCA) tracks the ratio of non-dominant, localized research perspectives versus standardized academic tropes across tasks.

\section{Analytical Pipeline and Reproducibility Protocol}

\subsection{Pipeline Steps}
Automated API Execution $\rightarrow$ Text Preprocessing $\rightarrow$ Vector Corpus Construction $\rightarrow$ High-Dimensional Clustering \& Statistical Variance Analysis.

\subsection{Validity and Reproducibility Safeguards}
Execution is constrained within tightly defined temporal windows to mitigate silent back-end model updates by proprietary providers, with full raw JSON archiving for complete auditability.
